\section{Next Greater Element I}
\label{sec:sq-next-greater-1}

\noindent We now turn to the second major group of problems in this
chapter: the \textbf{monotonic stack}, introduced in
Section~\ref{sec:sq-intro}. This first problem is the purest possible
instance of the pattern, and every later monotonic-stack problem in this
chapter is a variation or extension of the technique derived here.

\problemstatement{We are given two arrays, $\textit{nums1}$ and
$\textit{nums2}$, both containing distinct integers, where every element
of $\textit{nums1}$ is also present somewhere in $\textit{nums2}$. For
each element of $\textit{nums1}$, find its \textbf{next greater element}
in $\textit{nums2}$: the first element to its right, within
$\textit{nums2}$, that is strictly greater than it. If no such element
exists, the answer for that element is $-1$.}

\begin{patternbox}
\emph{``Next greater element''} is, almost word for word, the keyword
from Section~\ref{sec:sq-intro}'s pattern list. The structure here is
two-part: first solve ``next greater element, for every position, within
\textit{nums2}'' (a self-contained monotonic-stack problem), then a simple
hash-map lookup translates that solved array into answers for
$\textit{nums1}$.
\end{patternbox}

\subsection*{Understanding the problem carefully}

The core sub-problem is: for every index $i$ in $\textit{nums2}$, find the
nearest index $j>i$ with $\textit{nums2}[j] > \textit{nums2}[i]$, or
report none exists. A brute-force solution checks every pair, costing
$O(n^2)$. The monotonic stack from Section~\ref{sec:sq-intro} solves this
in $O(n)$ by scanning right to left and maintaining a stack that is always
strictly decreasing from bottom to top.

\begin{ideabox}[Monotonic Decreasing Stack, Scanned Right to Left]
Scan $\textit{nums2}$ from right to left. Before processing
$\textit{nums2}[i]$, pop every stack element $\le \textit{nums2}[i]$
(they can never be the answer for anything still to the left, since
$\textit{nums2}[i]$ is closer and at least as large). Whatever remains on
top of the stack after this popping (if anything) is exactly
$i$'s next greater element. Then push $\textit{nums2}[i]$.
\end{ideabox}

\subsection*{Why this gives the correct next greater element for every index}

After popping every $\le \textit{nums2}[i]$ value, the stack's remaining
top -- if any -- is the \emph{nearest} value to the right of $i$ that
exceeds $\textit{nums2}[i]$. This holds because the stack only ever
contains values in decreasing order from bottom to top, and every value
currently on the stack represents some index to the right of $i$ that has
not yet been ``defeated'' by anything closer and at least as large; the
topmost (most recently pushed, hence nearest) surviving value is therefore
guaranteed to be the closest qualifying greater value. Once computed for
every index of $\textit{nums2}$, store the results in a hash map keyed by
value (values are distinct, so this is unambiguous), and answer each
query for $\textit{nums1}$ with a single $O(1)$ lookup.

\subsection*{Algorithm}

\begin{enumerate}
  \item Initialize an empty stack and an empty hash map
        $\textit{nextGreater}$.
  \item For $i$ from $n-1$ down to $0$ (scanning $\textit{nums2}$ right
        to left):
  \begin{enumerate}
    \item While the stack is not empty and its top $\le
          \textit{nums2}[i]$: pop.
    \item $\textit{nextGreater}[\textit{nums2}[i]] \gets$ the stack's top
          if non-empty, else $-1$.
    \item Push $\textit{nums2}[i]$.
  \end{enumerate}
  \item For each query value in $\textit{nums1}$: look up
        $\textit{nextGreater}[\textit{value}]$.
\end{enumerate}

\begin{algorithm}[H]
\caption{Next Greater Element I}
\begin{algorithmic}[1]
\Procedure{NextGreaterElement}{\textit{nums1}, \textit{nums2}}
  \State $\textit{stack} \gets$ empty, $\textit{nextGreater} \gets$ empty map
  \For{$i \gets n-1$ \textbf{down to} $0$}
    \While{\textit{stack} not empty \textbf{and} top $\le \textit{nums2}[i]$}
      \State pop
    \EndWhile
    \State $\textit{nextGreater}[\textit{nums2}[i]] \gets$ (\textit{stack} empty $? -1 :$ top)
    \State push $\textit{nums2}[i]$
  \EndFor
  \State \Return $[\textit{nextGreater}[x]$ for each $x$ in $\textit{nums1}]$
\EndProcedure
\end{algorithmic}
\end{algorithm}

\subsection*{Dry run}

\dryrun{Take $\textit{nums2} = [1,3,4,2]$, $\textit{nums1} = [4,1,2]$.}

Scanning right to left:

\begin{itemize}
  \item $i=3$ (value $2$): stack empty. $\textit{nextGreater}[2]=-1$.
        Push $2$. Stack $=[2]$.
  \item $i=2$ (value $4$): top is $2 \le 4$, pop. Stack empty.
        $\textit{nextGreater}[4]=-1$. Push $4$. Stack $=[4]$.
  \item $i=1$ (value $3$): top is $4 > 3$, no pop.
        $\textit{nextGreater}[3]=4$. Push $3$. Stack $=[4,3]$.
  \item $i=0$ (value $1$): top is $3 > 1$, no pop.
        $\textit{nextGreater}[1]=3$. Push $1$. Stack $=[4,3,1]$.
\end{itemize}

Map: $\{2{:}{-1}, 4{:}{-1}, 3{:}4, 1{:}3\}$. Querying $\textit{nums1} =
[4,1,2]$: answers $[-1, 3, -1]$.

\subsection*{C++ implementation}

\begin{lstlisting}
#include <bits/stdc++.h>
using namespace std;

vector<int> nextGreaterElement(vector<int>& nums1, vector<int>& nums2) {
    unordered_map<int,int> nextGreater;
    stack<int> st;

    for (int i = nums2.size() - 1; i >= 0; i--) {
        while (!st.empty() && st.top() <= nums2[i]) {
            st.pop();
        }
        nextGreater[nums2[i]] = st.empty() ? -1 : st.top();
        st.push(nums2[i]);
    }

    vector<int> result;
    for (int x : nums1) result.push_back(nextGreater[x]);
    return result;
}

int main() {
    vector<int> nums2 = {1, 3, 4, 2};
    vector<int> nums1 = {4, 1, 2};
    for (int x : nextGreaterElement(nums1, nums2)) cout << x << " ";
    cout << endl;
    return 0;
}
\end{lstlisting}

\begin{complexitybox}
\textbf{Time Complexity.} Each element of $\textit{nums2}$ is pushed once
and popped at most once across the entire scan, giving $O(n)$ for the
preprocessing; each of the $m$ queries is an $O(1)$ map lookup, giving an
overall time complexity of
\[
  O(n + m).
\]
\textbf{Space Complexity.} $O(n)$ for the stack and hash map.
\end{complexitybox}

\begin{takeawaybox}
The monotonic stack derived here -- scan right to left, pop everything
$\le$ the current element, the survivor on top is the answer, then push
-- is the master template for this entire sub-chapter. Every remaining
monotonic-stack section either applies this template directly (next
smaller, previous greater/smaller) or builds one more layer on top of it
(histogram area, trapping rain water, subarray sums).
\end{takeawaybox}
